
%%%%%%%%%%%%%%%%%%%%%%%%%%%%%%%%%%%%%%%%
%           main body
%%%%%%%%%%%%%%%%%%%%%%%%%%%%%%%%%%%%%%%%

%-----------------------------------
%	TITLE PAGE
%-----------------------------------

\title[第九章\\
数项级数]{\texorpdfstring{\LARGE \bfseries 第九章\quad 数项级数}{第九章 数项级数}} % The short title appears at the bottom of every slide, the full title is only on the title page
\author[\smaller \S 9.2\\
上极限与下极限] % Your institution as it will appear on the bottom of every slide, maybe shorthand to save space
{\Large \bfseries \S 9.2 上极限与下极限}
% \author{Singular Value Decomposition} % Your name
% \date{\today} % Date, can be changed to a custom date
\date{}

%------------------------------------------------

\begin{document}

%------------------------------------------------

\begin{frame}

\titlepage % Print the title page as the first slide

\end{frame}

%-----------------------------------
%	PRESENTATION SLIDES
%-----------------------------------

\begin{frame}

\begin{center}
\begin{tikzpicture}[
scale=1.5,
x=0.7cm,y=1.2cm,
declare function={
    damped_wave(\x)=exp(-0.2*\x)*sin(deg(\x*3));  % 阻尼振动函数
    upper_env(\x)=exp(-0.2*\x);  % 上包络线
    lower_env(\x)=-exp(-0.2*\x);  % 下包络线
}
]
 
% 绘制坐标轴
\draw[->] (-0.5,0) -- (6.5,0) node[right] {$x$};
\draw[->] (0,-1.2) -- (0,1.5) node[above] {$y$};

% 绘制阻尼振动曲线（蓝色实线）
\draw[blue,thick,samples=100,smooth] 
    plot[domain=0:6] (\x,{damped_wave(\x)});
 
% 绘制包络线（红色虚线）
\draw[red,dashed,samples=50] 
    plot[domain=0:6] (\x,{upper_env(\x)})
    plot[domain=0:6] (\x,{lower_env(\x)});
\end{tikzpicture}
\end{center}

\vspace*{2em}

我们引入比 (序列) 极限更弱一些的概念, 即所谓的\alert{上、下极限}.

\end{frame}

%------------------------------------------------

\section[定义]{上、下极限的定义}

%------------------------------------------------

\begin{frame}{数列的极限点}

\begin{Def}
在有界数列 $\{x_n\}$ 中, 若存在一个子列 $\{x_{n_k}\},$ 使得 $\lim\limits_{k \to \infty} x_{n_k} = \xi,$ 则称 $\xi$ 为数列 $\{x_n\}$ 的一个\alert{极限点}.
\end{Def}

\pause

极限点的概念可以等价地定义为: $\xi$ 的任意邻域中都包含 $\{x_n\}$ 中的无穷多项.记
$$E = \{ \xi\in\mathbb{R} ~:~ \xi \text{ 是 $\{x_n\}$ 的极限点} \},$$
对于有界数列, 由于它总有收敛子列, 因此 $E$ 非空.

\begin{attnblock}
数列极限点的概念要与点集拓扑中极限点的概念相区别. 点集拓扑中, 一个点 $\xi$ 是集合 $S$ 的极限点, 若 $\xi$ 的任意去心邻域与 $S$ 的交非空.
\pause
二者的唯一区别在于, 当 $\xi$ 本身在 $\{x_n\}$ 中出现无穷多次时, $\xi$ 是 $\{x_n\}$ 的极限点, 但可能只是相应点集的孤立点.

数列极限点与点集拓扑中附着点 (adhere point) 的概念更接近.
\end{attnblock}

\end{frame}

%------------------------------------------------

\begin{frame}

\end{frame}

%------------------------------------------------

\section[运算]{上、下极限的运算}

%------------------------------------------------

\begin{frame}

\end{frame}

%------------------------------------------------

\end{document}
